\documentclass[11pt,a4paper]{article}
\usepackage[margin=27mm,top=25mm,bottom=25mm]{geometry}
\usepackage[T1]{fontenc}
\usepackage{amsmath,amsthm,mathtools}
\usepackage{newtxtext,newtxmath}
\usepackage{microtype}
\usepackage{booktabs,tabularx,array}
\usepackage{enumitem}
\usepackage{xcolor}
\usepackage[hidelinks]{hyperref}
\usepackage{fancyhdr,lastpage}
\usepackage{listings}
\usepackage{caption}
\usepackage{float}
\usepackage{needspace}
\definecolor{ink}{HTML}{1C2934}
\definecolor{dim}{HTML}{65727D}
\definecolor{accent}{HTML}{1C5265}
\hypersetup{pdftitle={The least trinomial multiple of a primitive degree-61 binary polynomial},pdfauthor={Sergey Nefedov},pdfsubject={Exact computation and a Lean proof of polynomial divisibility}}
\newtheorem{theorem}{Theorem}
\newtheorem{lemma}[theorem]{Lemma}
\newtheorem{proposition}[theorem]{Proposition}
\newtheorem{corollary}[theorem]{Corollary}
\newcommand{\F}{\mathbb F}
\newcommand{\Mod}{\mathop{\mathrm{mod}}}
\newcommand{\poly}{x^{61}+x^{45}+x^{32}+x^2+1}
\newcommand{\Bmax}{2^{31}}
\newcommand{\xor}{\mathbin{\mathrm{xor}}}
\pagestyle{fancy}
\fancyhf{}
\fancyhead[L]{\small\color{dim} S. Nefedov \;\textbar\; TheoremDB P2436}
\fancyhead[R]{\small\color{dim} October 2026}
\fancyfoot[C]{\small\color{dim} \thepage\ / \pageref{LastPage}}
\renewcommand{\headrulewidth}{0.3pt}
\setlength{\headheight}{14pt}
\setlength{\parindent}{1.25em}
\setlength{\parskip}{0.15em}
\setlist[itemize]{leftmargin=*,itemsep=2pt,topsep=4pt}
\lstset{basicstyle=\ttfamily\footnotesize,breaklines=true,frame=single,rulecolor=\color{dim!30},numbers=none,columns=fullflexible,xleftmargin=0.6em,xrightmargin=0.6em}
\begin{document}
\thispagestyle{empty}
\vspace*{8mm}
{\small\bfseries\color{accent} RESEARCH PREPRINT \hfill 17 OCTOBER 2026}\par
\vspace{5mm}
{\huge\bfseries\color{ink} The Least Trinomial Multiple of a Primitive\\[1pt] Degree-61 Binary Polynomial}\par
\vspace{4mm}
{\Large\color{dim} Exact exhaustive search and a Lean verification of the witnesses}\par
\vspace{9mm}
{\large\bfseries Sergey Nefedov}\par
{\normalsize Faculty of Mechanics and Mathematics, Lomonosov Moscow State University}\par
\vspace{1mm}
{\small Moscow, Russia}\par
\vspace{9mm}
\noindent\textbf{Abstract.}\;
We determine the least degree of a trinomial multiple of
\(f(x)=x^{61}+x^{45}+x^{32}+x^2+1\) in \(\F_2[x]\) within the range
\(0<a<b\le 2^{31}\). The least such degree is
\(b=848\,144\,186\), attained at \(a=301\,682\,885\).
An exhaustive, deterministic search over all \(2^{31}\) positive exponents
finds exactly three exponent pairs in the prescribed range.
The algorithm uses an exact, 60-bit parity-erased key for each element of a
primitive degree-61 finite field, followed by externally partitioned radix
sorting and a full recovery pass. We prove that collisions are equivalent to
trinomial divisibility and that the enumeration cannot skip any smaller
exponent. Independent polynomial remainder calculations verify the three
witnesses. A Lean~4 and Mathlib development formalizes the bitset computation
and its implication for divisibility in \((\mathbb Z/2\mathbb Z)[X]\); the
minimality claim is supported by exhaustive computation, not yet by a Lean
kernel proof. Source code and a reproducibility guide accompany this article.
\par\vspace{4mm}
\noindent\textbf{Keywords:} finite fields; sparse polynomial multiple;
trinomials; exhaustive computation; LFSR; Lean~4; computer-assisted proof.\par
\noindent\textbf{MSC 2020:} 12E20, 12Y05, 68W30, 11T71.\par
\vspace{7mm}\hrule\vspace{6mm}
\section{Problem and contribution}
Consider the monic polynomial
\begin{equation}\label{eq:f}
 f(x)=\poly\in\F_2[x].
\end{equation}
TheoremDB problem P2436 asks for the smallest integer \(b\le\Bmax\)
for which there exists \(a\) with \(0<a<b\) and
\begin{equation}\label{eq:divisibility}
 f(x)\mid 1+x^a+x^b.
\end{equation}
An earlier recorded exact computation excluded \(b\le 2^{28}\)
\cite{theoremdb}. That prefix is not a solution of the requested minimization
problem. Our contribution is an explicit least degree and a complete
algorithmic exclusion of every smaller degree, together with source code
that permits full external reproduction.

\begin{theorem}[Least trinomial multiple]\label{thm:main}
For \(f\) given in \eqref{eq:f}, the smallest integer \(b\) satisfying
\eqref{eq:divisibility} with \(0<a<b\le 2^{31}\) is
\[
 \boxed{b_0=848\,144\,186,\qquad a_0=301\,682\,885.}
\]
More precisely, the complete set of pairs \((a,b)\) in this range is
\begin{align*}
 &(301\,682\,885,\;848\,144\,186),\\
 &(452\,471\,585,\;1\,293\,170\,928),\\
 &(603\,365\,770,\;1\,696\,288\,372).
\end{align*}
\end{theorem}
The computational part of the proof is specified in
Sections~\ref{sec:algorithm}--\ref{sec:results} and released with this
preprint. In particular, \emph{the proof of minimality does not follow from
verifying the three positive witnesses alone}; it relies on the complete
collision enumeration.

\section{Exact finite-field model}
Set \(K=\F_2[x]/(f)\), and write \(\alpha=x\bmod f\). The recorded
field-model certificate and our standalone verifier check Rabin's
irreducibility criterion for the prime degree 61:
\begin{equation}\label{eq:rabin}
 x^{2^{61}}\equiv x\pmod f,\qquad
 \gcd(f,x^2+x)=1.
\end{equation}
The same verifier evaluates the 59-step Lucas--Lehmer sequence modulo
\(2^{61}-1\), obtaining zero. Therefore \(2^{61}-1\) is prime. Since
\(\alpha\ne0,1\), it follows that
\begin{equation}\label{eq:order}
 \operatorname{ord}(\alpha)=2^{61}-1.
\end{equation}
In particular, the residues \(\alpha^1,\ldots,\alpha^{2^{31}}\)
are pairwise distinct.

Represent an element of \(K\) by the integer whose 61 low binary digits
are its polynomial coefficients. Let
\[
 M=2^{61}-1,\qquad
 T=2^{45}+2^{32}+2^2+1.
\]
For a state \(s\in\{0,\ldots,M\}\), multiplication by \(\alpha\)
is represented exactly by
\begin{equation}\label{eq:transition}
 \operatorname{step}(s)=
 ((2s)\bmod 2^{61})\ \xor\
 \begin{cases}T,&\text{if bit 60 of }s\text{ is 1},\\0,&\text{otherwise}.\end{cases}
\end{equation}
Here \(\xor\) is bitwise exclusive OR; it implements addition in
characteristic two. Equation~\eqref{eq:transition} follows directly from
\(x^{61}\equiv x^{45}+x^{32}+x^2+1\pmod f\).
Starting from state \(1\), the \(i\)-th application produces the exact
representation \(s_i\) of \(\alpha^i\).

\begin{lemma}[Parity-erased collision criterion]\label{lem:collision}
For \(1\le i<j\le 2^{31}\), put \(q_i=\lfloor s_i/2\rfloor\).
Then
\begin{equation}\label{eq:equiv}
 q_i=q_j \quad\Longleftrightarrow\quad
 f(x)\mid x^j+x^i+1.
\end{equation}
\end{lemma}
\begin{proof}
The quotient by two discards exactly the constant coefficient. Equality
of two keys therefore says that \(s_i\) and \(s_j\) differ at most in
bit zero. Their residues are distinct by \eqref{eq:order}, so the constant
coefficients must differ, and \(s_i\xor s_j=1\). Conversely, if
\(s_i\xor s_j=1\), all higher bits agree, hence the keys agree.
Because XOR implements addition in \(K\), the former identity is
\(\alpha^i+\alpha^j+1=0\), equivalent to divisibility by \(f\).
\end{proof}

The use of distinctness is important: without it, a repeated \(s_i=s_j\)
would also give equal keys but would not give \(s_i\xor s_j=1\).
Here distinctness is guaranteed by the verified field order; candidate
verification additionally checks the exact XOR condition.

\section{Complete exact search}\label{sec:algorithm}
By Lemma~\ref{lem:collision}, the task is equivalent to enumerating all
duplicate values in the stream
\begin{equation}\label{eq:stream}
 q_1,q_2,\ldots,q_B,\qquad B=2^{31}.
\end{equation}
Each \(q_i\) fits into 60 bits. A key occupies an unsigned 64-bit word;
thus the uncompressed raw stream uses
\(8B=17\,179\,869\,184\) bytes (16~GiB). To avoid holding it in RAM,
the program partitions keys according to the top eight bits
\[
 p(q)=q\gg52\in\{0,1,\ldots,255\}.
\]
Every equal pair has equal partition index, so partitioning cannot
separate a genuine match.

\begin{enumerate}[leftmargin=2em,itemsep=4pt]
\item \textbf{Generate.} Iterate \eqref{eq:transition} exactly once for
every \(i=1,\ldots,B\), including the last endpoint. Append \(q_i\)
to file \(p(q_i)\); count each record and check file sizes.
\item \textbf{Partition-local exact sort.} Load each of 256 files in
turn, sort full 64-bit words using five stable 11-bit radix passes,
and compare adjacent words for literal equality. Since the top eight
bits are constant within a partition, the remaining low 52 bits are
fully ordered by those passes. Record each duplicate key.
\item \textbf{Recover exponents.} Reiterate \(s_1,\ldots,s_B\). For
each duplicate key recover both exponent labels and check that their
full states XOR to one and \(0<a<b\le B\).
\item \textbf{Check invariants.} Confirm that 256 partition counts sum
to \(B\), the byte counts agree, each collision key has exactly two
exponents, and the second pass ends at the same field state.
\end{enumerate}
No hash fingerprint or probabilistic filter is used to accept or reject
keys. The key truncation in Lemma~\ref{lem:collision} is exact for this
question, and all sorted comparisons use the entire stored key.

\begin{proposition}[Completeness of the search]\label{prop:complete}
If the algorithm above terminates with all 256 partitions successfully
processed and the stated invariants satisfied, its output is precisely
the set of pairs \((a,b)\) with \(0<a<b\le B\) and
\(f\mid x^b+x^a+1\).
\end{proposition}
\begin{proof}
Every \(i\) in the inclusive range \(1\le i\le B\) is visited by
the generation loop, with exactly one key and exactly one partition.
For each solution, Lemma~\ref{lem:collision} gives a pair of equal keys;
such keys lie in the same partition. Radix sorting compares exact keys;
all duplicates become adjacent, so the pair is detected. Since there
are only two 61-bit states with any fixed 60-bit key, and the powers
are distinct by \eqref{eq:order}, no key can occur more than twice.
The recovery pass assigns the unique two exponent labels and verifies
the defining XOR equality, excluding spurious candidates. Conversely,
any emitted pair passes precisely the equivalence in the lemma.
\end{proof}

The finite number of radix passes is independent of \(B\); for fixed
word size, generation, radix sorting and recovery require \(O(B)\)
word operations, although external I/O and memory-bandwidth costs are
substantial. Disk space is \(8B\) bytes for the keys plus a few
smaller metadata files. Processing a partition containing \(n_p\)
records requires \(16n_p\) bytes for the two radix buffers and a small
fixed counting array. The code releases each bucket after successful
verification. The raw files are regeneration inputs, not retained as
part of the distributable source archive.

\section{Results and exact arithmetic checks}\label{sec:results}
The complete computation reports the three exponent pairs listed in
Theorem~\ref{thm:main}. The table also lists the independently computed
remainders of the monomials modulo \(f\); within each row, the two
61-bit values differ exactly in bit zero.
\begin{table}[H]
\centering\small
\caption{Exact polynomial remainders for the three witnesses.}
\label{tab:witnesses}
\begin{tabular}{@{}r r l l@{}}\toprule
$a$ & $b$ & $x^a\bmod f$ & $x^b\bmod f$\\\midrule
301\,682\,885 & 848\,144\,186 & \texttt{116ffd86894d162e} & \texttt{116ffd86894d162f}\\
452\,471\,585 & 1\,293\,170\,928 & \texttt{0a3dd96b1419e23d} & \texttt{0a3dd96b1419e23c}\\
603\,365\,770 & 1\,696\,288\,372 & \texttt{0abd85a640cd46f8} & \texttt{0abd85a640cd46f9}\\\bottomrule
\end{tabular}
\end{table}
The verifier \texttt{verify/verify\_witness.py} computes these residues
by schoolbook carryless polynomial multiplication followed by generic
polynomial reduction, an implementation separate from the C++ LFSR
transition and radix search. It also compares the two arithmetic
methods on a prefix of exponents.

The third row has an algebraic explanation:
\begin{equation}\label{eq:square}
 (x^{848144186}+x^{301682885}+1)^2
 =x^{1696288372}+x^{603365770}+1
 \quad\text{in }\F_2[x].
\end{equation}
Because characteristic two eliminates cross terms, it is the square of
the first solution. The second row is a distinct solution, not that
squaring. The whole search includes even-even pairs, so this square is
not accidentally discarded.

\subsection{Machine run and consistency checks}
\IfFileExists{run_results.tex}{The reproducible C++17 implementation completed an inclusive scan of
\(B=2\,147\,483\,648\) positive exponents, successfully closing all
256 of 256 partitions and finding exactly three duplicate keys. A
second sweep recovered exactly the three exponent pairs of
Theorem~\ref{thm:main}. The process exited with code~0. Summary
statistics from its generated log are shown below.
\begin{center}
\small
\begin{tabular}{@{}l r@{}}\toprule
Quantity & Observed value\\\midrule
Enumerated states (inclusive) & 2\,147\,483\,648\\
Partition files processed & 256 of 256\\
Total equal-key pairs & 3\\
Final field state, decimal & 263\,596\,493\,844\,334\,377\\
Smallest / largest partition sizes & 8\,379\,809 / 8\,397\,622\\
Wall-clock time & 355.454 s\\
Peak resident set size (process) & 133\,468 KiB (130.34 MiB)\\
Temporary raw key volume & 16 GiB\\\bottomrule
\end{tabular}
\end{center}
The machine-readable summary and process log are included as
\texttt{results/full\_run\_summary.json} and
\texttt{results/full\_search.log}. Their totals and all three
witnesses pass the bundled independent Python metadata and remainder
checks. The 16~GiB raw intermediate is not shipped with the archive;
its exact regeneration is prescribed by the program. This run is a
computational certificate whose assurance depends on correct execution
of the published source, not a proof checked by the Lean kernel.
}{The full-run summary
will be inserted from the completed machine-readable execution log.}
The main program writes a JSON summary with the inclusive endpoint,
all 256 individual bucket counts, the duplicate keys, recovered
exponents, the final LFSR state, and elapsed time. A prefix test
matches the historical TheoremDB residue after 2,000,000 iterations:
\(2\,021\,773\,856\,510\,245\,769\), with no solutions in that
prefix. All three exhibited witnesses pass the independent polynomial
remainder checker.

\section{Lean proof artifacts and their scope}
The source distribution includes two Lean files, supplied in the
checked Lean package associated with this project:
\begin{itemize}
\item \texttt{P2436.lean}, a pure-Lean bitset program with theorems
\texttt{first\_candidate}, \texttt{second\_candidate}, and
\texttt{squared\_candidate}, using \texttt{decide +kernel} rather than
\texttt{native\_decide}.
\item \texttt{P2436Poly.lean}, a Mathlib development of the
bitset-to-polynomial correspondence in the quotient by~\(f\). Its
soundness theorem lifts a successful Boolean check to ordinary
polynomial divisibility, yielding three witness theorems.
\end{itemize}
The files are pinned to Lean~4.35.0 and Mathlib v4.35.0 in the archive.
Their contents were inspected for \texttt{sorry}, \texttt{admit}, and
custom \texttt{axiom} declarations; none are present.
\textbf{The present packaging environment does not contain Lean and
could not re-run \texttt{lake build}.} Therefore the Lean compilation
claims in the supplied archive remain to be verified on an independent
Lean installation or by the included GitHub Actions workflow.
Neither Lean file proves exhaustive minimality. The division witnesses
and exhaustive minimality are two different proof obligations.

\section{Reproduction and limitations}
The source release contains a standalone C++17 exhaustive search, an
independent Python polynomial-arithmetic verifier, the Rabin and
Lucas--Lehmer field checks, and regression tests. The essential commands
from the repository root are:
\begin{lstlisting}
sh tests/test_search.sh
c++ -O3 -std=c++17 src/exact_search.cpp -o exact_search
./exact_search 2147483648 experiments/full_data
python3 verify/verify_witness.py --summary experiments/full_data/run_summary.json
lake build
\end{lstlisting}
The full search requires at least 16~GiB of available disk space for
raw keys, plus headroom, and enough RAM for the largest two-buffer
partition. It is recommended to run it on a native 64-bit platform
with adequate free disk and stable power. The raw 16~GiB stream is not
included with the GitHub source ZIP; it can be reconstructed exactly.
The primary scan was executed in this study, but an independent
\emph{second full search program} is not included in this release.
Source audit, independent witness verification and externally
reproducing the full exhaustive search remain important for
high-assurance use of the minimality conclusion.

Prior algorithms for low-weight polynomial multiples, including
Zech-logarithm and discrete-logarithm methods, are studied by
Didier--Laigle-Chapuy \cite{didier} and
Peterlongo--Sala--Tinnirello \cite{peterlongo}; sparse multiples
in a broader algebraic setting are studied in \cite{giesbrecht}.
We do not claim that radix sorting or out-of-core collision
search is a new algorithmic paradigm. The result is the computed
least degree for this particular, previously documented instance.
A literature search identified no earlier public value for this
specific minimal pair, but this is not an exhaustive priority audit.

\section*{Code and data availability}
The complete source, exact enumeration record, and reproduction
instructions accompany this preprint as a GitHub-ready archive. The raw
temporary 16 GiB key stream is not included; deterministic regeneration
is described and implemented.\par

\clearpage
\begin{thebibliography}{9}
\bibitem{theoremdb}
TheoremDB, \emph{P2436: Least trinomial multiple of a primitive degree-61
polynomial}, research record, reviewed 24 July 2026.
\url{https://theoremdb.org/statements/P2436/}.

\bibitem{didier}
F. Didier and Y. Laigle-Chapuy, \emph{Finding low-weight polynomial
multiples using discrete logarithm}, IEEE International Symposium on
Information Theory, 2007, pp.~1036--1040.
\url{https://arxiv.org/abs/cs/0701069}.

\bibitem{peterlongo}
P. Peterlongo, M. Sala, and C. Tinnirello, \emph{A discrete logarithm-based
approach to compute low-weight multiples of binary polynomials},
Finite Fields and Their Applications \textbf{38} (2016), 57--71.
\url{https://doi.org/10.1016/j.ffa.2015.12.005}.

\bibitem{giesbrecht}
M. Giesbrecht, D. Roche, and H. Tilak, \emph{Computing sparse multiples
of polynomials}, arXiv:1009.3214.
\url{https://arxiv.org/abs/1009.3214}.
\end{thebibliography}

\section*{Appendix A. Reproducibility audit}
\noindent\textbf{Observed environment.} The fresh exhaustive search was
compiled with GCC 14.2.0 on Linux x86-64 (Intel Xeon Platinum 8573C).
A separate Python interpreter was used for the reference polynomial
remainder computations and the field-property checks. The program
reported an elapsed time of 355.454 seconds, a peak RSS of
133\,468 KiB, and successful termination of the complete scan.

\noindent\textbf{Recorded integrity checks.} The SHA-256 digests of the
completed machine-readable summary and process log are, respectively:
\begin{lstlisting}
9531cf943ab223532be089c498c9f54ee55187ba8b7d8f50039235df239f18bc
7c8ee6ced0def78db8fa84c8a180197808eba549dfc378e9eedb86960dca562d
\end{lstlisting}
These digests protect file integrity only. They are not substitutes for
the correctness proof in Proposition~\ref{prop:complete}, auditing the
C++ enumeration, or an independent full rerun.

\noindent\textbf{Concrete replay.} A minimal independent audit uses:
\begin{lstlisting}
python3 verify/verify_field.py
python3 verify/verify_witness.py
python3 verify/audit_summary.py results/full_run_summary.json
sh tests/test_search.sh
\end{lstlisting}
The complete scan command in Section~6 additionally regenerates all
256 partitions and verifies their exact keys before removing them. The
partial-reference Python dictionary implementation may be used for
smaller bounds, as a second method to check the indexing conventions.

\noindent\textbf{Limits of formal certification.} The supplied Lean
files are included unchanged from their checked-source package. A
source inspection found no proof holes, but a new Lean kernel replay
requires an external Lean environment. Even after all Lean witness
theorems compile, the universal minimality statement is still
established by the exhaustive C++ enumeration, not by a Lean theorem.
\end{document}
